\documentclass{beamer}
\usepackage[utf8]{inputenc}

\usetheme{Madrid}
\usecolortheme{default}
\usepackage{amsmath,amssymb,amsfonts,amsthm}
\usepackage{txfonts}
\usepackage{tkz-euclide}
\usepackage{listings}
\usepackage{adjustbox}
\usepackage{array}
\usepackage{tabularx}
\usepackage{gvv}
\usepackage{lmodern}
\usepackage{circuitikz}
\usepackage{tikz}
\usepackage{graphicx}
\usepackage{multicol}

\setbeamertemplate{page number in head/foot}[totalframenumber]

\usepackage{tcolorbox}
\tcbuselibrary{minted,breakable,xparse,skins}

\definecolor{bg}{gray}{0.95}
\DeclareTCBListing{mintedbox}{O{}m!O{}}{%
  breakable=true,
  listing engine=minted,
  listing only,
  minted language=#2,
  minted style=default,
  minted options={%
    linenos,
    gobble=0,
    breaklines=true,
    breakafter=,,
    fontsize=\small,
    numbersep=8pt,
    #1},
  boxsep=0pt,
  left skip=0pt,
  right skip=0pt,
  left=25pt,
  right=0pt,
  top=3pt,
  bottom=3pt,
  arc=5pt,
  leftrule=0pt,
  rightrule=0pt,
  bottomrule=2pt,
  toprule=2pt,
  colback=bg,
  colframe=orange!70,
  enhanced,
  overlay={%
    \begin{tcbclipinterior}
    \fill[orange!20!white] (frame.south west) rectangle ([xshift=20pt]frame.north west);
    \end{tcbclipinterior}},
  #3,
}
\lstset{
    language=C,
    basicstyle=\ttfamily\small,
    keywordstyle=\color{blue},
    stringstyle=\color{orange},
    commentstyle=\color{green!60!black},
    numbers=left,
    numberstyle=\tiny\color{gray},
    breaklines=true,
    showstringspaces=false,
}

\numberwithin{equation}{section}
\lstset{
  language=Python,
  basicstyle=\ttfamily\small,
  keywordstyle=\color{blue},
  stringstyle=\color{orange},
  numbers=left,
  numberstyle=\tiny\color{gray},
  breaklines=true,
  showstringspaces=false
}


\title{Problem 9.4.2}
\author{Sarvesh Tamgade}

\date{\today} 
\begin{document}

\begin{frame}
\titlepage
\end{frame}

\section{Question}
\begin{frame}{Question}
\textbf{Question}:
Find the roots of the following quadratic equation graphically:
 \begin{align}
  x\brak{2x + 3} = x^2 + 1
 \end{align}


\end{frame}

\section{Solution}
\begin{frame}[fragile]
    \frametitle{Solution}
First, we expand and rearrange the equation to standard form:
\begin{align}
x\brak{2x + 3} &= x^2 + 1\\
2x^2 + 3x &= x^2 + 1\\
2x^2 + 3x - x^2 - 1 &= 0\\
x^2 + 3x - 1 &= 0
\end{align}

\textbf{Input Variables:  }\\
The given quadratic can be written in the conic form:
\begin{align}
\vec{x}^T\vec{V}\vec{x} + 2\vec{u}^T\vec{x} + f = 0
\end{align}

where
\begin{align}
\vec{V} = \myvec{1 & 0 \\ 0 & 0}, \quad \vec{u} = \myvec{\frac{3}{2} \\ 0}, \quad f = -1
\end{align}

\end{frame}
\begin{frame}[fragile]
    \frametitle{Solution}
Since the roots of the quadratic correspond to intersections with the x-axis, we represent the line:
\begin{align}
L: \vec{x} = \vec{h} + \kappa\vec{m}
\end{align}

with
\begin{align}
\vec{h} = \myvec{0 \\ 0}, \quad \vec{m} = \myvec{1 \\ 0}
\end{align}

\end{frame}
\begin{frame}[fragile]
    \frametitle{Solution}

\begin{table}[H]
\centering
\begin{tabular}{|c|c|}
\hline
\textbf{Symbol} & \textbf{Value} \\
\hline
$\vec{V}$ & $\myvec{1 & 0 \\ 0 & 0}$ \\
\hline
$\vec{u}$ & $\myvec{\frac{3}{2} \\ 0}$ \\
\hline
$f$ & $-1$ \\
\hline
$\vec{h}$ & $\myvec{0 \\ 0}$ \\
\hline
$\vec{m}$ & $\myvec{1 \\ 0}$ \\
\hline
\end{tabular}
\caption{Input Parameters}
\label{tab:input}
\end{table}

\end{frame}

\begin{frame}[fragile]
    \frametitle{Solution}
The points of intersection of a line with a conic are given by:
\begin{align}
\kappa = \frac{1}{\vec{m}^T\vec{V}\vec{m}} \sbrak{-\vec{m}^T\brak{\vec{V}\vec{h} + \vec{u}} \pm \sqrt{\brak{\vec{m}^T\brak{\vec{V}\vec{h} + \vec{u}}}^2 - g\brak{\vec{h}}\brak{\vec{m}^T\vec{V}\vec{m}}}}
\end{align}

where
\begin{align}
g\brak{\vec{h}} = \vec{h}^T\vec{V}\vec{h} + 2\vec{u}^T\vec{h} + f
\end{align}

\textbf{Step 1:} Compute $\vec{m}^T\vec{V}\vec{m}$
\begin{align}
\vec{m}^T\vec{V}\vec{m} = \myvec{1 & 0}\myvec{1 & 0 \\ 0 & 0}\myvec{1 \\ 0} = 1
\end{align}

\textbf{Step 2:} Compute $\vec{V}\vec{h} + \vec{u}$
\begin{align}
\vec{V}\vec{h} + \vec{u} = \myvec{1 & 0 \\ 0 & 0}\myvec{0 \\ 0} + \myvec{\frac{3}{2} \\ 0} = \myvec{\frac{3}{2} \\ 0}
\end{align}

\end{frame}

\begin{frame}[fragile]
    \frametitle{Solution}
\textbf{Step 3:} Compute $\vec{m}^T\brak{\vec{V}\vec{h} + \vec{u}}$
\begin{align}
\vec{m}^T\brak{\vec{V}\vec{h} + \vec{u}} = \myvec{1 & 0}\myvec{\frac{3}{2} \\ 0} = \frac{3}{2}
\end{align}

\textbf{Step 4:} Compute $g\brak{\vec{h}}$
\begin{align}
g\brak{\vec{h}} = \vec{h}^T\vec{V}\vec{h} + 2\vec{u}^T\vec{h} + f = 0 + 0 + \brak{-1} = -1
\end{align}

\textbf{Step 5:} Substitute into formula for $\kappa$
\begin{align}
\kappa &= -\frac{3}{2} \pm \sqrt{\brak{\frac{3}{2}}^2 - \brak{-1}\brak{1}}\\
&= -\frac{3}{2} \pm \sqrt{\frac{9}{4} + 1}\\
&= \frac{-3 \pm \sqrt{13}}{2}
\end{align}

\end{frame}
\begin{frame}[fragile]
    \frametitle{Solution}
\textbf{Step 6:} Evaluate roots
\begin{align}
\kappa_1 &= \frac{-3 + \sqrt{13}}{2}\\
\kappa_2 &= \frac{-3 - \sqrt{13}}{2}
\end{align}

\textbf{Step 7:} Find intersection points
The intersection points are obtained as:
\begin{align}
\vec{x}_1 &= \vec{h} + \kappa_1\vec{m} = \myvec{0 \\ 0} + \frac{-3 + \sqrt{13}}{2}\myvec{1 \\ 0} = \myvec{\frac{-3 + \sqrt{13}}{2} \\ 0}\\
\vec{x}_2 &= \vec{h} + \kappa_2\vec{m} = \myvec{0 \\ 0} + \frac{-3 - \sqrt{13}}{2}\myvec{1 \\ 0} = \myvec{\frac{-3 - \sqrt{13}}{2} \\ 0}
\end{align}

\end{frame}
\begin{frame}[fragile]
    \frametitle{Solution}
Thus, the quadratic $x^2 + 3x - 1 = 0$ intersects the x-axis at:
\begin{align}
x &= \frac{-3 - \sqrt{13}}{2} \approx -3.303\\
x &= \frac{-3 + \sqrt{13}}{2} \approx 0.303
\end{align}

The roots of the original equation $x\brak{2x + 3} = x^2 + 1$ are:
\begin{align}
\boxed{x = \frac{-3 \pm \sqrt{13}}{2}}
\end{align}
\end{frame}


\section{Graph}
\begin{frame}
    \frametitle{Graph}
    \begin{figure}[htbp]
    \centering
    \includegraphics[width=0.65\linewidth]{FIG/fig1.png}
    \caption{Vector Representation}
    \label{fig:FIG/fig1.png}
\end{figure}
\end{frame}
\section{ C Code}
\begin{frame}[fragile]
\frametitle{C Code }
\begin{lstlisting}[language=C]
#include <math.h>

void intersect_line_conic(
    const double V[4],   
    const double u[2],     
    double f,              
    const double h[2],     
    const double m[2],     
    double kappa_out[2],   
    double x1[2],
    double x2[2],
    int *status            
) {
    double Vm0 = V[0]*m[0] + V[1]*m[1];
    double Vm1 = V[2]*m[0] + V[3]*m[1];
    double mTVm = m[0]*Vm0 + m[1]*Vm1;
    
\end{lstlisting}
\end{frame}
\begin{frame}[fragile]
\frametitle{C Code }
\begin{lstlisting}[language=C]
 double Vh0 = V[0]*h[0] + V[1]*h[1];
    double Vh1 = V[2]*h[0] + V[3]*h[1];
    
    double Vh_u0 = Vh0 + u[0];
    double Vh_u1 = Vh1 + u[1];
    
    double mT_Vh_u = m[0]*Vh_u0 + m[1]*Vh_u1;
    
    double hTVh = h[0]*Vh0 + h[1]*Vh1;
    double two_uTh = 2.0*(u[0]*h[0] + u[1]*h[1]);
    double g = hTVh + two_uTh + f;
    
    double disc = mT_Vh_u*mT_Vh_u - g*mTVm;
\end{lstlisting}
\end{frame}
\begin{frame}[fragile]
\frametitle{C Code }
\begin{lstlisting}[language=C]
   if (disc > 1e-12) {
        double r = sqrt(disc);
        kappa_out[0] = (-mT_Vh_u + r) / mTVm;
        kappa_out[1] = (-mT_Vh_u - r) / mTVm;
        
        x1[0] = h[0] + kappa_out[0]*m[0]; x1[1] = h[1] + kappa_out[0]*m[1];
        x2[0] = h[0] + kappa_out[1]*m[0]; x2[1] = h[1] + kappa_out[1]*m[1];
        
        *status = 2;
\end{lstlisting}
\end{frame}
\begin{frame}[fragile]
\frametitle{C Code }
\begin{lstlisting}[language=C]
    } else if (fabs(disc) <= 1e-12) {
        kappa_out[0] = kappa_out[1] = (-mT_Vh_u) / mTVm;
        x1[0] = x2[0] = h[0] + kappa_out[0]*m[0];
        x1[1] = x2[1] = h[1] + kappa_out[0]*m[1];
        
        *status = 1;
    } else {
        *status = 0;
    }
}
\end{lstlisting}
\end{frame}
\begin{frame}[fragile]
\frametitle{Python + C Code }
\begin{lstlisting}[language=Python]
import ctypes as ct
import numpy as np
import matplotlib.pyplot as plt

# --- load the shared library (same folder) ---
lib = ct.CDLL("./libsimple_conic.so")

lib.intersect_line_conic.argtypes = [
    ct.POINTER(ct.c_double),  # V[4]
    ct.POINTER(ct.c_double),  # u[2]
    ct.c_double,              # f
    ct.POINTER(ct.c_double),  # h[2]
    ct.POINTER(ct.c_double),  # m[2]
    ct.POINTER(ct.c_double),  # kappa[2] out
    ct.POINTER(ct.c_double),  # x1[2] out
    ct.POINTER(ct.c_double),  # x2[2] out
    ct.POINTER(ct.c_int)      # status out
]
\end{lstlisting}

\end{frame}
\begin{frame}[fragile]
    \frametitle{Python + C Code}
    \begin{lstlisting}
lib.intersect_line_conic.restype = None

# ---- helper: convert quadratic ax^2+bx+c into conic form ----
def quadratic_to_conic(a, b, c):
    V = np.array([a, 0.0, 0.0, 0.0], dtype=np.double)  # [[a,0],[0,0]]
    u = np.array([b/2.0, 0.0], dtype=np.double)        # so 2u^T x = b x
    f = np.double(c)
    return V, u, f

# ---- our quadratic: x^2 + 3x - 1 = 0 ----
a, b, c = 1.0, 3.0, -1.0
V, u, f = quadratic_to_conic(a, b, c)

# Intersect with x-axis (y=0): line x = h + k m
h = np.array([0.0, 0.0], dtype=np.double)
m = np.array([1.0, 0.0], dtype=np.double)

\end{lstlisting}
\end{frame}

\begin{frame}[fragile]
    \frametitle{Python + C Code }
    \begin{lstlisting}
# Outputs (allocated for C)
kappa = np.zeros(2, dtype=np.double)
x1 = np.zeros(2, dtype=np.double)
x2 = np.zeros(2, dtype=np.double)
status = ct.c_int(0)

# ---- call the C function ----
lib.intersect_line_conic(
    V.ctypes.data_as(ct.POINTER(ct.c_double)),
    u.ctypes.data_as(ct.POINTER(ct.c_double)),
    ct.c_double(f),
    h.ctypes.data_as(ct.POINTER(ct.c_double)),
    m.ctypes.data_as(ct.POINTER(ct.c_double)),
    kappa.ctypes.data_as(ct.POINTER(ct.c_double)),
    x1.ctypes.data_as(ct.POINTER(ct.c_double)),
    x2.ctypes.data_as(ct.POINTER(ct.c_double)),
    ct.byref(status)
)

\end{lstlisting}
\end{frame}
\begin{frame}[fragile]
    \frametitle{Python + C Code }
    \begin{lstlisting}

# Collect results
roots = []
if status.value >= 1:
    roots.append(float(x1[0]))
    if status.value == 2:
        roots.append(float(x2[0]))
    roots.sort()

print("status:", status.value)
print("roots:", roots)  # expected [(-3-√13)/2, (-3+√13)/2] ≈ [-3.303, 0.303]

# ---- plot parabola and roots ----
xs = np.linspace(-5, 2, 600)  # adjusted range for our equation
ys = a*xs*xs + b*xs + c

\end{lstlisting}
\end{frame}
\begin{frame}[fragile]
    \frametitle{Python + C Code }
    \begin{lstlisting}
plt.figure()
plt.plot(xs, ys, label=rf"$y={a:.0f}x^2{b:+.0f}x{c:+.0f}$")
plt.axhline(0, linestyle="--", linewidth=1, label="x-axis")

for r in roots:
    plt.scatter([r], [0.0], s=60, zorder=3, label=f"root {r:.3f}")

plt.xlabel("x"); plt.ylabel("y")
plt.title("Roots via line–conic intersection (NumPy + C)")
plt.grid(True)
plt.legend()
plt.savefig("parabola.png")
plt.show()
\end{lstlisting}
\end{frame}

\begin{frame}[fragile]
\frametitle{Python Code }
\begin{lstlisting}[language=Python]
import numpy as np
import matplotlib.pyplot as plt

def quadratic_to_conic(a, b, c):
    V = np.array([[a, 0.0],
                  [0.0, 0.0]], dtype=float)
    u = np.array([b/2.0, 0.0], dtype=float)  # so that 2 u^T x = b x
    f = float(c)
    return V, u, f

def line_conic_intersection(V, u, f, h, m, eps=1e-12):
    # m^T V m
    Vm = V @ m
    mTVm = float(m @ Vm)
    
    # Vh + u
    Vh_u = V @ h + u
\end{lstlisting}

\end{frame}
\begin{frame}[fragile]
    \frametitle{Python Code}
    \begin{lstlisting}

    mT_Vh_u = float(m @ Vh_u)
    
    # g(h) = h^T V h + 2 u^T h + f
    g = float(h @ (V @ h) + 2.0 * (u @ h) + f)
    
    disc = mT_Vh_u**2 - g * mTVm
    
    if disc > eps:
        r = np.sqrt(disc)
        k1 = (-mT_Vh_u + r) / mTVm
        k2 = (-mT_Vh_u - r) / mTVm
        
        X1 = h + k1 * m
        X2 = h + k2 * m
        
        return 2, np.array([k1, k2], dtype=float), np.vstack([X1, X2])

\end{lstlisting}
\end{frame}

\begin{frame}[fragile]
    \frametitle{Python Code}
    \begin{lstlisting}
    elif abs(disc) <= eps:
        k = (-mT_Vh_u) / mTVm
        X = h + k * m
        
        return 1, np.array([k, k], dtype=float), np.vstack([X, X])
    else:
        return 0, np.array([np.nan, np.nan]), np.array([[np.nan, np.nan],
                                                         [np.nan, np.nan]])

# ---------- Problem setup ----------
# Given quadratic: x^2 + 3x - 1 = 0
a, b, c = 1.0, 3.0, -1.0

# Conic parameters (V, u, f)
V, u, f = quadratic_to_conic(a, b, c)

\end{lstlisting}
\end{frame}
\begin{frame}[fragile]
    \frametitle{Python Code}
    \begin{lstlisting}
# x-axis as the line: y = 0 -> h = (0,0), m = (1,0)
h = np.array([0.0, 0.0], dtype=float)
m = np.array([1.0, 0.0], dtype=float)

# ---------- Solve via line–conic intersection ----------
status, kappa, X = line_conic_intersection(V, u, f, h, m)

# Roots are the x-coordinates of intersection points with y=0
roots = []
if status >= 1:
    roots = sorted([float(X[0, 0]), float(X[1, 0])]) if status == 2 else [float(X[0, 0])]

print("Status (2:two real, 1:tangent, 0:none):", status)
print("kappa values:", kappa)
print("Intersection points (x,y):\\n", X)
print("Roots (x-intercepts):", roots)

\end{lstlisting}
\end{frame}
\begin{frame}[fragile]
    \frametitle{Python Code}
    \begin{lstlisting}
# ---------- Plot ----------
xs = np.linspace(-5, 2, 600)
ys = a*xs**2 + b*xs + c

plt.figure()
plt.plot(xs, ys, label=rf"$y={a:.0f}x^2{b:+.0f}x{c:+.0f}$")
plt.axhline(0, linestyle="--", linewidth=1, label="x-axis")


\end{lstlisting}
\end{frame}
\begin{frame}[fragile]
    \frametitle{Python Code}
    \begin{lstlisting}
if status >= 1:
    plt.scatter([X[0,0]], [0.0], s=60, zorder=3, label=f"root {X[0,0]:.3f}")
    if status == 2:
        plt.scatter([X[1,0]], [0.0], s=60, zorder=3, label=f"root {X[1,0]:.3f}")

plt.xlabel("x"); plt.ylabel("y")
plt.title("Roots via line–conic intersection (vectors & matrices)")
plt.grid(True)
plt.legend()
plt.savefig("fig1.png")
plt.show()
\end{lstlisting}
\end{frame}
\end{document}
